%\chapter{Integration and Transcendental Functions}

\documentclass{ximera}

\graphicspath{ 
{./}
{./chapters/ch1Preliminary/}
}
\definecolor{programcode}{gray}{0.65}
\newenvironment{programcode}[1]{\ignorespaces
%\def\stmtopen##1{##1}%
%\formtmp{programcode}{#1}
}{\noindent\textcolor{programcode}{\rule{\columnwidth}{1pt}}\vskip2pt\par\addvspace{\baselineskip}}%
\renewcommand{\caption}[1]{\text{#1}}
%\newcounter{figure}
%\renewcommand{\thefigure}{\arabic{figure}}
%
%\renewcommand{\caption}[1]{%
%  \refstepcounter{figure}%
%  \textbf{Figure \thefigure:} #1%
%}
\author{Leon Kaganovskiy}
\title{Integration and Transcendental Functions}
\begin{document}
\begin{abstract}
This chapter introduces integration and its applications. We use SageMath to explore definite and indefinite integrals numerically, symbolically, and graphically, helping us understand the relationship between integration and differentiation. We apply these ideas to calculate areas, volumes, and other quantities that arise in mathematics, science, and practical problems.
\end{abstract}
\maketitle

%  \begin{image}[15cm] %%\sidecaption[t]   remove ch2 etc reference local directory

\section{Basic Integration}

SageMath can compute rather complex integrals requiring substantial analytical work.  Given that our goal in this book is not just to give SageMath commands but rather use it to illustrate Mathematical techniques, we compute many Calculus integrals even though they are just "done" in SageMath a\-la "black box computation" without many analytical techniques taught in Calculus 2.  In the cell below, we start with basic integrations with power rule ($\int x^n dx = \frac{x^{n+1}}{n+1} \ \ n \ne -1$) as well as trigonometric functions.
$$ \int x^3 dx = \frac{x^4}{4} \ \ \ \ \ \ \ \ \ \ \ \ \ \ \ \ \int \sqrt{x} dx = \int x^{1/2} dx = \frac{x^{3/2}}{\frac{3}{2}} =\frac{2x^{3/2}}{3} $$

\begin{programcode}{SageMath Code and Output:}
\begin{verbatim}
f(x) = x^3;  F(x) = integrate(f(x),x)
show('f(x) = ',f(x),'  its antiderivative F(x) = ',F(x))
f(x) = sqrt(x);  F(x) = integrate(f(x),x)
show('f(x) = ',f(x),'  its antiderivative F(x) = ',F(x))
f(x) = cos(x);  F(x) = integrate(f(x),x)
show('f(x) = ',f(x),'  its antiderivative F(x) = ',F(x))
f(x) = sin(x);  F(x) = integrate(f(x),x)
show('f(x) = ',f(x),'  its antiderivative F(x) = ',F(x))
f(x) = (sec(x))^2;  F(x) = integrate(f(x),x)
show('f(x) = ',f(x),'  its antiderivative F(x) = ',F(x))
\end{verbatim}
\end{programcode}
The output is shown in Fig \ref{IntPowerTrig}
\begin{image}[15cm]
%\sidecaption[t]
\includegraphics[scale=.65]{/IntPowerTrig.png}
\caption{Basic integrals of power and trig functions}
\label{IntPowerTrig}       % Give a unique label
\end{image}

The next example illustrates the simplification of monomial division resulting in a straightforward power integration, but SageMath does not need explicit simplification.

\begin{programcode}{SageMath Code and Output:}
\begin{verbatim}
f(x) = (x^3-3*x^2+2*x-1)/sqrt(x);  F(x) = integrate(f(x),x);  
show('f(x) = ',f(x),'  simplified f(x) = ', expand(f(x)) )
show('  its antiderivative F(x) = ',F(x))
\end{verbatim}
\end{programcode}
The output is shown in Fig \ref{MonomialDivision}
\begin{image}[15cm]
%\sidecaption[t]
\includegraphics[scale=.65]{/MonomialDivision.png}
\caption{Monomial division simplification}
\label{MonomialDivision}       % Give a unique label
\end{image}


The indefinite integral is determined up to a constant.  In the problem below, an extra condition is specified to determine that constant.
$$ f(x) = x^3 \ \ \ \ \Rightarrow \ \ \ \ F(x) = \int x^3 dx = \frac{x^4}{4} + C $$
$$ F(1) =  \frac{1}{4} + C =2 \ \ \ \ \Rightarrow   C = \frac{7}{4}  \ \ \ \ \Rightarrow  F(x) =  \frac{x^4}{4} + \frac{7}{4}$$
We also show that the graphs of various antiderivatives are just vertical shifts of each other.

\begin{programcode}{SageMath Code and Output:}
\begin{verbatim}
f(x) = x^3;  F(x) = integrate(f(x),x);  x1 = 1;  F1 = 2
show('f(x) = ',f(x))
show('  its antiderivative F(x) = ',F(x))
C1 = F1 - F(1); 
show(' F(1) = ', F(1), '   C1 = ',C1 )
show('   particular solution = F(x) + C1= ', F(x)+C1)
plot([F(x),F(x)-2,F(x)-1,F(x)+2,F(x)+1], (x,-2,2))
\end{verbatim}
\end{programcode}
The output is shown in Fig \ref{AntiderivGraph}
\begin{image}[15cm]
%\sidecaption[t]
\includegraphics[scale=.65]{/AntiderivGraph.png}
\caption{Antiderivatives graph}
\label{AntiderivGraph}       % Give a unique label
\end{image}

The SageMath code below illustrates a classical Physics problem of a ball thrown up into the air.  Assume that the initial velocity is 48 ft/sec and the initial height is 24 feet.   Using Newton's 2nd law, we can describe its trajectory and determine when it reaches the max height and what this height is.
$$ a=-32 ft/sec^2 \ \ \ \ \ \Rightarrow v(t)=\int a dt = \int (-32)dt = -32t+C \ \ \ \ \ \Rightarrow$$
$$ v(0) = C = v_0 = 48  \ \ \ \ \ \Rightarrow v(t)= -32t+48 \ \ \ \ \ \Rightarrow$$
$$ h(t)= \int v(t)dt = \int  (-32t+48) dt = -32 \frac{t^2}{2} +48t + C= -16t^2+48t + C  \ \ \ \ \ \Rightarrow$$
$$ h(0)=C=h_0=24  \ \ \ \ \ \Rightarrow h(t) =-16t^2+48t +24 $$
Maximum height is reached when velocity is 0:
$$v(t)= -32t+48 = 0  \ \ \ \ \ \Rightarrow t_m = \frac{48}{32} = 1.5 \ \ \ \ \ \Rightarrow$$
$$ h_{max} = h(t_m)= -16\cdot 1.5^2+48 \cdot 1.5 +24 = 60 $$

\begin{programcode}{SageMath Code and Output:}
\begin{verbatim}
var('t');  a = -32;  v0 = 48;  h0 = 24;
v(t) = integrate(a,t) + v0;
h(t) = integrate(v(t),t) + h0;
show('velocity = v(t) = ', v(t))
show('height = h(t) = ', h(t))
tm = -v0/a;  
show('Max height of h(tm) = ',h(tm))
show('    is reached at time tm = ',tm)
a = 0; b = 4; c = -80; d = 65; # plot bounds 
plot1 = plot(v(t),(t,a,b), ymin=c, ymax=d, \
  legend_label= "v(t)", color='red', gridlines=True)
plot2 = plot(h(t),(t,a,b), ymin=c, ymax=d, \
  legend_label= "h(t)", color='blue', gridlines=True)
plot_all = plot1 + plot2; plot_all.show()

\end{verbatim}
\end{programcode}
The output is shown in Fig \ref{BallTrajectory}
\begin{image}[15cm]
%\sidecaption[t]
\includegraphics[scale=.65]{/BallTrajectory.png}
\caption{Vertical projectile trajectory}
\label{BallTrajectory}       % Give a unique label
\end{image}

The SageMath code below illustrates left and right Riemann sums converging to definite integral (area under the function).  In addition, SageMath can symbolically compute the sum and find the limit as the number of subdivisions $n \rightarrow \infty$.

\begin{programcode}{SageMath Code and Output:}
\begin{verbatim}
f(x)=x^2;  a = 0; b = 2
n = 10;  delx = N((b-a)/n);  sumL = 0;  sumR = 0
for j in range(1,n+1):
    sumL = sumL + f(a + (j-1)*delx)*delx
    sumR = sumR + f(a +     j*delx)*delx
    
print('For n = ',n, 'subdivisions: ')
print('   sumL = ', sumL, '  sumR = ', sumR)

n = 1000;  delx = N((b-a)/n);  sumL = 0;  sumR = 0
for j in range(1,n+1):
    sumL = sumL + f(a + (j-1)*delx)*delx
    sumR = sumR + f(a +     j*delx)*delx
    
print('For n = ',n, 'subdivisions: ')
print('   sumL = ', sumL, '  sumR = ', sumR)

var('n,k');   delx = (b-a)/n
sumLn = sum(f(a + (k-1)*delx)*delx,k,1,n);
show('sumLn = ', sumLn)
IntL = limit(sumLn,n=infinity)
print('sumLn converges to definite integral (area) = IntL = ', IntL)
Int = integrate(f(x),x,a,b)
print('Integrate command gives = ', Int)
p = plot(f(x), (x,a,b), fill=True) + plot(f(x), (x,a,b)); show(p)
\end{verbatim}
\end{programcode}
The output is shown in Fig \ref{RiemannSum}
\begin{image}[15cm]
%\sidecaption[t]
\includegraphics[scale=.65]{/RiemannSum.png}
\caption{Riemann sum estimation}
\label{RiemannSum}       % Give a unique label
\end{image}



In the problem below, SageMath is used to illustrate the Fundamental Theorem of Calculus and how positive and negative areas between the function and $x$-axis add up to the definite integral.  
$$ f(x) = \frac{x^2-x-6}{\sqrt{x}} = x^{\frac{3}{2}}- x^{\frac{1}{2}}-6x^{-\frac{1}{2}}  \Rightarrow $$
$$\ \ \ \ \ F(x)=\int f(x)dx = \frac{x^{\frac{5}{2}}}{\frac{5}{2}}- \frac{x^{\frac{3}{2}}}{\frac{3}{2}}-6\frac{x^{\frac{1}{2}}}{\frac{1}{2}} =  \frac{2x^{\frac{5}{2}}}{5}- \frac{2x^{\frac{3}{2}}}{3}-12x^{\frac{1}{2}}$$
Using the calculations below
$$ \int_1^4 f(x)dx = F(4)-F(1) = -\frac{64}{15} = A_1 (area \ below) + A_2 (area \ above) =  $$
$$ \left( -\frac{52\sqrt{3}}{5}+ \frac{184}{15} \right) +   \left( \frac{52\sqrt{3}}{5}- \frac{248}{15} \right) $$

\begin{programcode}{SageMath Code and Output:}
\begin{verbatim}
a = 1; b = 4
f(x) = (x^2-x-6)/sqrt(x);  F(x) = integrate(f(x),x);  
show('f(x) = ',f(x),'  simplified f(x) = ', expand(f(x)))
show('  its antiderivative F(x) = ',F(x))

Int = integrate(f(x),x,a,b)
show('Integrate command gives = ', Int)
show('  From Fundamental Theorem F(b)-F(a) = ',F(b)-F(a))

sol = solve(f(x)==0,x); show('solving f(x)==0 gives:', sol)
show('In this case, there are negative and positive areas')
A1 = integrate(f(x),x,a,3);   A2 = integrate(f(x),x,3,b);
show('Negative area A1 = ', A1, ' positive area A2 = ', A2)
show('  they add up to = A1+A2 = ', A1+A2)
p = plot(f(x), (x,a,b), fill=True) + plot(f(x), (x,a,b)) + \
    text("A1",(1.5,-2)) + text("A2",(3.6,0.7)); show(p)
\end{verbatim}
\end{programcode}
The output is shown in Fig \ref{PositNegatArea}
\begin{image}[15cm]
%\sidecaption[t]
\includegraphics[scale=.65]{/PositNegatArea.png}
\caption{Integral as area}
\label{PositNegatArea}       % Give a unique label
\end{image}


In the SageMath code below we illustrate the computation of the average value.  
$$ f(x) = 4x^2-2x  \Rightarrow \ \ \ \ \ F(x)=\int f(x)dx = 4\frac{x^3}{3}- 2\frac{x^2}{2} =4\frac{x^3}{3}-x^2 $$
Using the calculations below
$$ \int_1^4 f(x)dx = F(4)-F(1) =69  \Rightarrow \ \ \ \ f_{ave}= \frac{1}{b-a}\int_a^b f(x)dx = \frac{69}{4-1}=23 $$The area under the graph  $\int_a^b f(x)dx = f_{ave} (b-a)$.  Therefore, the areas $A_1$ and $A_2$ shown in the graph above and below $y=f_{ave}$ are the same and the total area between the function $f(x)$ and the x-axis is equal to the area of the rectangle below $y=f_{ave}$.

\begin{programcode}{SageMath Code and Output:}
\begin{verbatim}
a = 1; b = 4
f(x) = 4*x^2-2*x;  F(x) = integrate(f(x),x);  
show('f(x) = ',f(x),'  simplified f(x) = ', expand(f(x)) )
show('  its antiderivative F(x) = ',F(x))

Int = integrate(f(x),x,a,b)
show('Integrate command gives = ', Int)
show('  also from Fundamental Theorem F(b)-F(a) = ',F(b)-F(a))
f_ave = Int/(b-a);  
show('Function average on [a,b] = f_ave = ',f_ave)

p = plot(f(x), (x,a,b), figsize = (4,3), fill=True) + plot(f(x), (x,a,b)) + \
    plot(f_ave,(x,a,b),color="red",thickness=3) + \
    plot(f_ave,(x,a,b), fill=True) + \
    text("A1",(1.3,15)) + text("A2",(3.6,32)); show(p)
\end{verbatim}
\end{programcode}
The output is shown in Fig \ref{AverageInt}
\begin{image}[15cm]
%\sidecaption[t]
\includegraphics[scale=.65]{/AverageInt.png}
\caption{Average of function}
\label{AverageInt}       % Give a unique label
\end{image}


Another illustration of the net change is provided by the velocity $v(t)$ of a particle moving along a straight line.  Its net change in position  $s(t)$ (forward for positive velocity and backward for negative velocity) is called displacement.
However, for the total distance, the regions of negative velocity (going backward - reducing displacement) have to be counted as positive, so the absolute value of velocity must be integrated.
$$ v(t) = s'(t) \Rightarrow \ \ \ \ \ Displacement = \int_a^b v(t)dt = s(b)-s(a)  $$
$$ Total \ distance = \int_a^b |v(t)|dt  $$
SageMath can integrate the absolute value of a function directly (with some warnings) as shown below.   However, a more illuminating approach is to identify regions of positive and negative velocity and integrate $-v(t) \ \ where \ \ v(t)<0$ as shown below as well.

\begin{programcode}{SageMath Code and Output:}
\begin{verbatim}
var('t'); a = 1;  b = 6;
v(t) = expand((t-1)*(t-4)*(t-6))
show('velocity = v(t) = ',v(t))
displacement = integrate(v(t),t,a,b)
show('The displacement = integrate v(t) = ', displacement)
total_dist = integrate(abs(v(t)),t,a,b)
d1 = integrate( v(t),t,1,4)
d2 = integrate(-v(t),t,4,6)  # note -v(t) sign
show('Total distance = integrate |v(t)| = ', total_dist)
show(  'we also check d1 + d2 = ', d1+d2)
plot(v(t),(t,a,b))
\end{verbatim}
\end{programcode}
The output is shown in Fig \ref{DisplaceDistance}
\begin{image}[15cm]
%\sidecaption[t]
\includegraphics[scale=.65]{/DisplaceDistance.png}
\caption{Distance and displacement}
\label{DisplaceDistance}       % Give a unique label
\end{image}


SageMath can easily evaluate integrals that require substitution.

\begin{programcode}{SageMath Code and Output:}
\begin{verbatim}
f(x) = x/sqrt(1-4*x^2);  F(x) = integrate(f(x),x)
show('f(x) = ',f(x),'  its antiderivative F(x) = ',F(x))
f(x) = x^3*cos(x^4+7);  F(x) = integrate(f(x),x)
show('f(x) = ',f(x),'  its antiderivative F(x) = ',F(x))
f(x) = x/(5*x^2-7)^3;  F(x) = integrate(f(x),x)
show('f(x) = ',f(x),'  its antiderivative F(x) = ',F(x))
\end{verbatim}
\end{programcode}
The output is shown in Fig \ref{ChangeVarIntegrat}
\begin{image}[15cm]
%\sidecaption[t]
\includegraphics[scale=.65]{/ChangeVarIntegrat.png}
\caption{Substitution}
\label{ChangeVarIntegrat}       % Give a unique label
\end{image}


It can confirm even and odd function integrals:

\begin{programcode}{SageMath Code and Output:}
\begin{verbatim}
f(x) = (sin(x))^5*(cos(x))^3;  F(x) = integrate(f(x),x)
show('f(x) = ',f(x),'  its antiderivative F(x) = ',F(x))
show('Integral on this odd function on symmetric limits: ',
      integrate(f(x),x,-pi/2,pi/2)) 
f(x) = (cos(x))^3;  F(x) = integrate(f(x),x)
show('f(x) = ',f(x),'  its antiderivative F(x) = ',F(x))
show('Integral on this even function on symmetric limits: ',
      integrate(f(x),x,-pi/2,pi/2)) 
show('  same as twice the integral on positive half: ',
      2*integrate(f(x),x,0,pi/2))
\end{verbatim}
\end{programcode}
The output is shown in Fig \ref{EvenOddInteg}
\begin{image}[15cm]
%\sidecaption[t]
\includegraphics[scale=.65]{/EvenOddInteg.png}
\caption{Even and odd functions}
\label{EvenOddInteg}       % Give a unique label
\end{image}



%%%%%%%%%%%%%%%%%%%%%%%%%%%%%%%%%%%%%%%%%%%%%%%%%%%%%%%%%%%%%%%%%%%%%%%
\section{Transcendental Functions Differentiation and Integration}

Before continuing with more specialized integration techniques, transcendental functions like exponential, logarithmic, and inverse trigonometric functions must be introduced.  There are different ways to introduce $\ln(x)$ and its derivative $(\ln(x))' = \frac{1}{x}$ in Calculus.  Some textbooks start with the exponential function $e^x$ and its derivative $(e^x)' = e^x$, introducing $\ln(x)$ as the inverse.  Others, start by defining $\ln(x) = \int_1^x \frac{1}{t}dt$ and use Second Fundamental of Calculus to derive $\frac{d}{dx}\ln(x) = \frac{d}{dx}\int_1^x \frac{1}{t}dt =\frac{1}{x}$.  Then, $e^x$ is introduced as the inverse of the logarithm.   In either case, these are standard derivatives that are used in Calculus.  Other bases are handled via replacing $b^x = e^{\ln(b)x}$ and change of log base formula $\log_b(x) = \frac{\ln(x)}{\ln(b)}$.    The log of the old base $\ln(b)$ is just a constant, so differentiation and integration are straightforward.  To summarize:
$$ (e^x)' =e^x \ \ \ \ \ (\ln(x))' = \frac{1}{x}  \ \ \ \ \  b^x = b^x\ln(b) \ \ \ \ \ \ (\log_b(x))' = \frac{1}{x\ln(b)}$$
$$  \int e^x dx = e^x \ \ \ \ \ \  \int \frac{1}{x}dx = \ln(x)  \ \ \ \ \ \int b^xdx = \frac{b^x}{\ln(b)}$$



\subsection{Logarithmic and exponential derivatives} %%%%%%%%%%%%%%%%%%%%%%%%%%%%%%%%%%%%%%%%%%%

The SageMath code below illustrates applications of these derivatives together with the chain rule.   Note also that even though in Calculus $\ln(x)=\log_e(x)$ and $\log(x) = \log_{10}(x)$, in SageMath and most other software (as well as in many advanced textbooks), $\log(x)$ is used to denote natural logarithm and a special notation like $\log(x,10)$ is used to denote logarithm base 10.  Any other base can be introduced by just replacing 10.  Note that as the output shows clearly, it is just immediately saved as a logarithm base $e$ via the base change formula.

\begin{programcode}{SageMath Code and Output:}
\begin{verbatim}
f(x) = log(x^3-3*x^2+7*x-4);  fp(x) = derivative(f(x),x)
show('f(x) = ',f(x)); show('  its derivative fp(x) = ',fp(x))
f(x) = log(x^3-3*x^2+7*x-4,10);  fp(x) = derivative(f(x),x)
show('f(x) = ',f(x)); show('  its derivative fp(x) = ',fp(x))
f(x) = e^(x^3-3*x^2+7*x-4);  fp(x) = derivative(f(x),x)
show('f(x) = ',f(x)); show('  its derivative fp(x) = ',fp(x))
f(x) = (log(4*x^2-3*x+2))^4;  fp(x) = derivative(f(x),x)
show('f(x) = ',f(x)); show('  its derivative fp(x) = ',fp(x))
f(x) = log(cos(x));  fp(x) = derivative(f(x),x)
show('f(x) = ',f(x)); show('  its derivative fp(x) = ',fp(x))
\end{verbatim}
\end{programcode}
The output is shown in Fig \ref{LogDerivat}
\begin{image}[15cm]
%\sidecaption[t]
\includegraphics[scale=.65]{/LogDerivat.png}
\caption{Basic log derivatives}
\label{LogDerivat}       % Give a unique label
\end{image}


Logarithms allow to simplify expressions involving multiplication/division and exponents before the differentiation.   For example:
$$ y =f(x) = \ln \left( \frac{x^5(x^3-1)^2}{\sqrt{7x^4+1}}  \right) = \ln(x^5(x^3-1)^2) - \ln((7x^4+1)^{1/2}) = $$
$$ \ln(x^5) + \ln((x^3-1)^2) - \ln((7x^4+1)^{1/2}) = $$
$$ 5\ln(x) + 2\ln(x^3-1) -\frac{1}{2} \ln(7x^4+1) \ \ \Rightarrow$$
$$ y' = f'(x) = \frac{5}{x} + 2\frac{1}{x^3-1}3x^2 - \frac{1}{2}\frac{1}{7x^4+1}28x^3 = $$
$$ \frac{5}{x} + \frac{6x^2}{x^3-1} - \frac{14x^3}{7x^4+1}$$
The SageMath can take such derivatives automatically, but the answer is not simplified and after the application of the full\_simplfy\(\), the analytical form is not the same.

There is a similar approach of logarithmic differentiation where no logarithms are involved to start with, but we apply it on both sides to considerably simplify the calculation.
$$ y =f(x) =\frac{x^5(x^3-1)^2}{\sqrt{7x^4+1}} \Rightarrow \ \ \ \ \ \ \ $$
$$ \ln(y) = \ln \left( \frac{x^5(x^3-1)^2}{\sqrt{7x^4+1}}  \right) =  5\ln(x) + 2\ln(x^3-1) -\frac{1}{2} \ln(7x^4+1) \ \ \Rightarrow$$
$$ \frac{1}{y}y' =  \frac{5}{x} + \frac{6x^2}{x^3-1} - \frac{14x^3}{7x^4+1} \ \ \ \ \Rightarrow $$
$$ y' =  \left( \frac{x^5(x^3-1)^2}{\sqrt{7x^4+1}}  \right) \left( \frac{5}{x} + \frac{6x^2}{x^3-1} - \frac{14x^3}{7x^4+1} \right)$$

\begin{programcode}{SageMath Code and Output:}
\begin{verbatim}
f(x) = log(x^5*(x^3-1)^2/sqrt(7*x^4+1));  
fp(x) = derivative(f(x),x)
show('f(x) = ',f(x))
show('  its derivative fp(x) = ',fp(x))
show(' apply full_simplify = ', fp(x).full_simplify())
f(x) = x^5*(x^3-1)^2/sqrt(7*x^4+1);  
fp(x) = derivative(f(x),x)
show('f(x) = ',f(x))
show(' its derivative fp(x) = ',fp(x))
show(' apply full_simplify = ', fp(x).full_simplify())
\end{verbatim}
\end{programcode}
The output is shown in Fig \ref{SimplifyLogDeriv}
\begin{image}[15cm]
%\sidecaption[t]
\includegraphics[scale=.65]{/SimplifyLogDeriv.png}
\caption{Log simplified derivatives}
\label{SimplifyLogDeriv}       % Give a unique label
\end{image}


Logarithms also allow to take derivatives of expressions with $x$ in both base and exponent:
$$ y = x^x \ \ \Rightarrow \ \ \ln(y) = \ln(x^x) = x\ln(x) \ \ \Rightarrow \ \ $$
$$ \frac{1}{y}y'= 1\cdot \ln(x) + x\frac{1}{x}  \ \ \  \ \ \Rightarrow \ \  y'= x^x(\ln(x) + 1)$$
SageMath has no trouble computing such derivatives either.

\begin{programcode}{SageMath Code and Output:}
\begin{verbatim}
f(x) = x^x;  fp(x) = derivative(f(x),x)
show('f(x) = ',f(x),'  its derivative fp(x) = ',fp(x))
\end{verbatim}
\end{programcode}
The output is shown in Fig \ref{XtoXderivat}
\begin{image}[15cm]
%\sidecaption[t]
\includegraphics[scale=.65]{/XtoXderivat.png}
\caption{$x^x$ derivative}
\label{XtoXderivat}       % Give a unique label
\end{image}




In the next two problems, we apply regular graphical analysis with first and second derivatives, min/max, inflection points etc... to logarithmic and exponential functions.  I repeat the code for the same function which was developed in the previous chapter.   First for the logarithm:
$$ f(x) = log(x^2+4x+7) \ \Rightarrow \ \ \ \ \ f'(x) = \frac{2x+4}{x^2+4x+7} \ \ \ $$
$$ f'(x) = 0 \ only \ at \ x=-2 \ \ \ \ \  x^2+4x+7 >0$$
$$ \ for \ any \ x \ (D=4^2-4\cdot 1 \cdot 7 = -12<0)$$

\begin{programcode}{SageMath Code and Output:}
\begin{verbatim}
def graphing_func(x,f,L,R,D,U):
    fp(x) = derivative(f(x),x); 
    fpp(x) = derivative(fp(x),x)
    show('f(x) = ', f(x))
    show('  derivative = fp(x) = ', fp(x))
    show('  factored form is =', factor(fp(x)) )
    show('  2nd derivative = fpp(x) = ', fpp(x))
    show('  factored form is =', factor(fpp(x)) )
    solp  = solve(fp(x)==0 ,x); 
    show('Zeros of derivative are: ', solp)
    solpp = solve(fpp(x)==0,x); 
    show('Zeros of 2nd derivative are: ', solpp)
    plot1 = plot(f(x)  ,(x,L,R), ymin=D, ymax=U, figsize = (5,4), \
        legend_label= "f(x)",color='red', gridlines=True)
    plot2 = plot(fp(x) ,(x,L,R), ymin=D, ymax=U, \
        legend_label= "fp(x)",color='blue', gridlines=True)
    plot3 = plot(fpp(x),(x,L,R), ymin=D, ymax=U, \
        legend_label= "fpp(x)",color='green', gridlines=True)
    plot_all = plot1 + plot2 + plot3; plot_all.show()
    
f(x) = log(x^2+4*x+7); 
graphing_func(x,f,-8,5,-1,4)
\end{verbatim}
\end{programcode}
The output is shown in Fig \ref{GraphLog}
\begin{image}[15cm]
%\sidecaption[t]
\includegraphics[scale=.65]{/GraphLog.png}
\caption{Graphing log function}
\label{GraphLog}       % Give a unique label
\end{image}


Next, consider the exponential appearing in a standard normal distribution.
$$ f(x) = e^{-x^2/2} \ \Rightarrow \ \ \ \ \ f'(x) = e^{-x^2/2} \frac{-2x}{2} =e^{-x^2/2} (-x) \   \ \ \text{therefore} $$
$$ f''(x) =e^{-x^2/2}(-x)(-x) + e^{-x^2/2}(-1)= e^{-x^2/2}(x^2-1) =  e^{-x^2/2}(x-1)(x+1) $$
It has the absolute maximum at $x=0$ and inflection points at $x=\pm 1$ marking concavity switching (positive/negative/positive) to produce a characteristic bell\-shape.  

\begin{programcode}{SageMath Code and Output:}
\begin{verbatim}
f(x) = e^(-x^2/2);   
graphing_func(x,f,-5,5,-1,1)
\end{verbatim}
\end{programcode}
The output is shown in Fig \ref{GraphExpX2}
\begin{image}[15cm]
%\sidecaption[t]
\includegraphics[scale=.65]{/GraphExpX2.png}
\caption{Normal exponential}
\label{GraphExpX2}       % Give a unique label
\end{image}


\subsection{Exponential models}%%%%%%%%%%%%%%%%%%%%%%%%%%%%%%%%%%%%%%%%%%%%%%

In the problem below, we use SageMath to illustrate Moore's law of doubling the number of transistors on computer chips.   We find the rate of change (derivative) and use \texttt{find\_root} to solve for the time when this growth rate reaches a given value.  Note that the base-2 doubling model is identical to the base-e model (the graphs overlap exactly).   The growth is so large that \texttt{plot\_semilogy()} must be used.
$$ n_1(t) = n_0 2^{t/d}, \ \ \ \ d = doubling \ time \ \ \ \ \ \ \ \ $$
$$ n_2(t) = n_0 e^{rt}, \ \ \ \ r = relative \ growth \ rate \ \ \ \ \ r = \frac{\ln(2)}{d}$$
$$ n_1'(t) = n_0 2^{t/d} \ln(2)\frac{1}{d} \ \ \ \ \  \ \  \ \  \ \ \ \  \ \ \  n_2'(t) = n_0 e^{rt}r$$
In addition, SageMath computes the average growth rate over the first 10 years:  $\frac{1}{10}\int_1^{10} n'(t) dt$.

\begin{programcode}{SageMath Code and Output:}
\begin{verbatim}
var('t')
n0 = 16 # number of transistors on the 1st chip
d = 18/12; # 18 months doubling time / 12 months;
r = N(log(2)/d);  
show('n0 = ', n0, ' doubling time = d = ',d)
show('   exp growth rate = r = ', r)
n1(t) = n0*2^(t/d);   n2(t) = n0*e^(r*t);  
np1(t) = derivative(n1(t),t);  
np2(t) = derivative(n2(t),t);  # rates (derivatives)
show('Exponential growth model base 2 = n(t) = ',n1(t))
show(' its rate of growth = np(t) = ', np1(t) )
show('Exponential growth model base e = n(t) = ',n2(t))
show('  its rate of growth = np(t) = ', np2(t))
rate1 = 100;   t1 = find_root(np1(t)==rate1,0,50); 
show('Rate reaches = np1 = ', rate1, ' at t1 = ', t1)
np_ave = 1/10*integral(np2(t),t,0,10); 
show('The average growth rate over 10 years = np_ave = ', np_ave)

a = 0; b = 50; c = 1; d = 20*10^9; # plot bounds
p1 = plot_semilogy(n1(t)  ,(t,a,b), ymin=c, 
  ymax=d, legend_label= "n1(t)",color='red', gridlines=True)
p2 = plot_semilogy(n2(t)  ,(t,a,b), ymin=c, 
  ymax=d, legend_label= "n2(t)",color='blue', gridlines=True)
a = 0; b = 50; c = 1; d = 200; # # plot bounds
p3 = plot(np2(t) ,(t,a,b), ymin=c, ymax=d, 
  legend_label= "np1(x)",color='blue', gridlines=True)
p4 = plot(rate1   ,(t,a,b), ymin=c, ymax=d, 
  legend_label= "n'=100",color='green', gridlines=True)
G = graphics_array([[p1+p2, p3+p4]])
G.show(figsize=[6, 3])
\end{verbatim}
\end{programcode}
The output is shown in Fig \ref{MooreLaw}
\begin{image}[15cm]
%\sidecaption[t]
\includegraphics[scale=.65]{/MooreLaw.png}
\caption{Moore Law}
\label{MooreLaw}       % Give a unique label
\end{image}


Let's consider car value depreciation as an example of exponential decay. Assume it reduces by 10\%  every year.  
 Then the value model and the rate of change are:
$$ n(t) = n_0(1-d)^t \ \ \ d =\text{decimal discount} $$
$$ \text{the rate of change:} \ \ n'(t) =  n_0(1-d)^t \ln(1-d)  $$
The corresponding base$-e$ model can be obtained as follows:
$$ n(t) = n_0 e^{-rt}  \ \Rightarrow \ \ \ \ n_0(1-d) = n_0 e^{-r\cdot 1}  \ \Rightarrow $$
$$ r = -\ln(1-d) \ \ \ \ \ \ \text{the rate of change: \ \ } n'(t) =   n_0 e^{-rt}(-r)$$
Analogous to the previous problem, SageMath is used to find the time when a particular rate is achieved, and the 
average rate over the first 10 years.

\begin{programcode}{SageMath Code and Output:}
\begin{verbatim}
n0 = 25000 # initial value of the car
d = 0.1; # discount percentage 10%;
r = N(-log(1-d));  T = 50;

show('n0 = ', n0, ' discount = d = ',d)
show('   exp rate = r = ', r)
n1(t) = n0*(1-d)^t;   n2(t) = n0*e^(-r*t);  
np1(t) = derivative(n1(t),t);  
np2(t) = derivative(n2(t),t);    # rates (derivatives)
show('Exponential decay model = n(t) = ',n1(t))
show('  its rate  = np(t) = ', np1(t) )
show('Exponential decay model base e = n(t) = ',n2(t))
show('  its rate  = np(t) = ', np2(t) )
rate1 = -1000;   t1 = find_root(np1(t)==rate1,0,50); 
show('The rate reaches value = np1 = ', rate1)
show('   at time t1 = ', t1)
np_ave = 1/10*integral(np2(t),t,0,10); 
show('Average rate of first 10 years = np_ave = ', np_ave)

p1 = plot(n1(t)  ,(t,0,T),  legend_label= "n1(t)",\
              color='red', gridlines=True)
p2 = plot(n2(t)  ,(t,0,T),  legend_label= "n2(t)", \
              color='blue', gridlines=True)
p3 = plot(np2(t) ,(t,0,T),  legend_label= "np1(t)", \
              color='blue', gridlines=True)
p4 = plot(rate1   ,(t,0,T), legend_label= "rate1", \
              color='green', gridlines=True)
G = graphics_array([[p1+p2, p3+p4]])
G.show(figsize=[6, 3])
\end{verbatim}
\end{programcode}
The output is shown in Fig \ref{CarValueExpDecay}
\begin{image}[15cm]
%\sidecaption[t]
\includegraphics[scale=.65]{/CarValueExpDecay.png}
\caption{Exponential decay}
\label{CarValueExpDecay}       % Give a unique label
\end{image}


A similar consideration can be applied to compound interest.
$$ \text{Discrete compounding:} \ F_n(t) = P \left( 1+ \frac{r}{n} \right)^{nt} \ \ \ \ \ \ $$
$$ \text{its rate of change:} \ F_n'(t) = P \left( 1+ \frac{r}{n} \right)^{nt}  \ln \left( 1+ \frac{r}{n} \right) n  $$
To find the time when it reaches a particular value:
$$F(t) = P \left( 1+ \frac{r}{n} \right)^{nt} = F_1 \ \Rightarrow \  \left( 1+ \frac{r}{n} \right)^{nt} = \frac{F_1}{P}  \ \Rightarrow \  $$
$$ \ln \left( \left( 1+ \frac{r}{n} \right)^{nt} \right) = \ln\left( \frac{F_1}{P} \right)    \ \Rightarrow  \  nt \ln \left( 1+ \frac{r}{n} \right) =  \ln\left( \frac{F_1}{P} \right) \Rightarrow \ \ \ \  $$
$$ t = \frac{ \ln \left( \frac{F_1}{P} \right)  }{ n \ln \left( 1+ \frac{r}{n} \right) }  $$
Analogously, for the continuous compounding:
$$ F_e(t) = Pe^{rt} \ \ \ \ \ \ F'(t) = Pe^{rt}r \Rightarrow \ \ \ \ \ \ \ F_e(t) = Pe^{rt} =F_1  \Rightarrow \ \ \ \ \ \  $$
$$ e^{rt} =\frac{F_1}{P}  \Rightarrow \ \ \ \ \ rt =  \ln \left( \frac{F_1}{P} \right)   \Rightarrow \ \ \ \ \ t = \frac{1}{r} \ln \left( \frac{F_1}{P} \right)  $$The SageMath code below computes and plots semiannual and continuous compoundings and their rates as well as times required to reach a particular future value.   The continuous compounding grows a bit faster.

\begin{programcode}{SageMath Code and Output:}
\begin{verbatim}
reset(); var('t')
P = 10000 # initial investment 
r = 0.10; # annual growth;
n = 2;  # number of compounding per year
F1 = 30000 # Future goal
T = 40 # max time for plots
Fn(t) = P*(1+r/n)^(n*t); # discrete compounding  
Fe(t) = P*e^(r*t);  # continous compounding
Fnp(t) = derivative(Fn(t),t);  # rate of change
Fep(t) = derivative(Fe(t),t);  # rate of change
show('Discrete compounding n =', n, ' Fn(t) =', Fn(t))
show('  its rate of growth = Fnp(t) = ', Fnp(t) )
show('Continuous compounding Fe(t) =', Fe(t))
show(' its rate of growth = Fep(t) = ', Fep(t) )
t1 = N(log(F1/P)/(n*log(1+r/n)));   t2 = N(log(F1/P)/r);
show('The time it takes Fn and Fe to reach F1 = ', F1)
show('   t1 = ', t1, ' and t2 = ', t2, ' respectively')

p1 = plot(Fn(t)  ,(t,0,T),  legend_label= "Fn(t)",\
              color='red', gridlines=True)
p2 = plot(Fe(t)  ,(t,0,T),  legend_label= "Fe(t)", \
              color='blue', gridlines=True)
p3 = plot(Fnp(t) ,(t,0,T),  legend_label= "Fnp(t)", \
              color='blue', gridlines=True)
p4 = plot(F1   ,(t,0,T), legend_label= "F1", \
              color='green', gridlines=True)
G = graphics_array([[p1+p2, p3+p4]])
G.show(figsize=[6, 3])
\end{verbatim}
\end{programcode}
The output is shown in Fig \ref{Compounding}
\begin{image}[15cm]
%\sidecaption[t]
\includegraphics[scale=.65]{/Compounding.png}
\caption{Compounding}
\label{Compounding}       % Give a unique label
\end{image}


The example below illustrates a logistic growth model.  A lake initially has 500 fish and the population increases logistically.  Find population at t = 0,1,10, infinity.   When is the inflection point?   How long does it take fish population to reach a particular value?   The formulas are somewhat tedious and SageMath helps enormously in the computations.   The computation of the inflection point requires a rather tedious second derivative, but SageMath finds it immediately.

\begin{programcode}{SageMath Code and Output:}
\begin{verbatim}
reset();   var('t');  T = 50 # max time for plots
F(t) = 10000/(1+19*e^(-t/5)); # logistic growth 
Finf = limit(F(t),t=infinity) 
show('Initial value F(0) = ', F(0))
show('Long term behaviour (carrying capacity) = ',Finf)
show('The population at time t=0: is =', N(F(0)))
show('The population at time t=10: is = ', N(F(10)))
Fp(t) = derivative(F(t),t);    
Fpp(t) = derivative(Fp(t),t);   # 1st (rate) and 2nd derivatives
show('Logistic model = F(t) = ', F(t))
show('  its rate of growth = Fp(t) = ', Fp(t) )

p1 = 0.3;  pF1 = N(Finf*p1);
t1 = find_root(F(t) == pF1,0,100);
show('The time it takes to reach p1 = ', p1*100,'% of ')
show('    carrying capacity is t1 = ', t1)
tpp = find_root(Fpp(t) == 0,0,100);  
show('The inflection point is at tpp = ',tpp)

p1 = plot(F(t)  ,(t,0,T), legend_label= "Fn(t)",\
              color='red', gridlines=True)
p2 = plot(pF1 ,(t,0,T), legend_label= "Finf*p1", \
              color='blue', gridlines=True)
p3 = plot(Fp(t) ,(t,0,T), legend_label= "Fnp(t)", \
              color='blue', gridlines=True)
p4 = plot(Fpp(t) ,(t,0,T), legend_label= "F1", \
              color='green', gridlines=True)
G = graphics_array([[p1+p2, p3+p4]])
G.show(figsize=[6, 3])
\end{verbatim}
\end{programcode}
The output is shown in Fig \ref{Logistic}
\begin{image}[15cm]
%\sidecaption[t]
\includegraphics[scale=.65]{/Logistic.png}
\caption{Logistic model}
\label{Logistic}       % Give a unique label
\end{image}


\subsection{Integrals with exponential and logarithmic functions} %%%%%%%%%%%%%%%%%%%%%%%%%%%%%%

As we mentioned before, SageMath has no trouble computing integrals involving change of variables, and integrals producing logarithms and exponentials are no exception.
$$ \int \frac{x}{x^2+1}dx = [u=x^2+1 \Rightarrow \ \ du=2xdx] =\int \frac{\frac{du}{2}}{u}=\frac{1}{2} \ln(u)=\frac{1}{2} \ln(x^2+1)$$
$$ \int_0^3 \frac{x}{x^2+1}dx = \int_1^{10} \frac{\frac{du}{2}}{u} = \frac{1}{2} \ln(u)|_1^{10}=\frac{1}{2}(\ln(10)-ln(1))=\frac{1}{2}\ln(10)$$
$$ \int \frac{2e^{4x}}{e^{4x}+7}dx = [u=e^{4x}+7 \Rightarrow \ \ du=e^{4x}4dx] =$$
$$ 2\int \frac{\frac{du}{4}}{u}=\frac{1}{2} \ln(u)=\frac{1}{2} \ln(e^{4x}+7)$$
$$ \int \tan(x)dx= \int \frac{\sin(x)}{\cos(x)}dx = [u=\cos(x) \Rightarrow \ \ du=-\sin(x)dx] =$$
$$ \int \frac{-du}{u}=- \ln(u)=- \ln|\cos(x)|$$
The answers generated by SageMath below are sometimes given in a different form than the theoretical ones above.

\begin{programcode}{SageMath Code and Output:}
\begin{verbatim}
a = 0; b = 3
f(x) = x/(x^2+1);  F(x) = integrate(f(x),x);  
show('f(x) = ',f(x), '  its antiderivative F(x) = ',F(x))

Int = integrate(f(x),x,a,b)
show('Integrate command gives = ', Int)
show('  Fundamental Theorem F(b)-F(a) = ',F(b)-F(a))

f(x) = 2*e^(4*x)/(e^(4*x)+7);  F(x) = integrate(f(x),x);  
show('f(x) = ',f(x), '  its antiderivative F(x) = ',F(x))

f(x) = tan(x);  F(x) = integrate(f(x),x)
show('f(x) = ',f(x),'  its antiderivative F(x) = ',F(x))
f(x) = sec(x);  F(x) = integrate(f(x),x)
show('f(x) = ',f(x),'  its antiderivative F(x) = ',F(x))
\end{verbatim}
\end{programcode}
The output is shown in Fig \ref{LogAntideriv}
\begin{image}[15cm]
%\sidecaption[t]
\includegraphics[scale=.65]{/LogAntideriv.png}
\caption{Antiderivatives involving logarithms}
\label{LogAntideriv}       % Give a unique label
\end{image}



The next example follows the previous one with $f(x) =  \frac{x}{x^2+1}$.  The highest power on the numerator below is the same as in the denominator $x^2$, so long division must be used to make a proper rational function, but here it is quite easy and SageMath computes it easily too.
$$\int \frac{x^2+x+1}{x^2+1} dx =\int 1+\frac{x}{x^2+1} dx = x + \frac{1}{2}\ln(x^2+1)$$

\begin{programcode}{SageMath Code and Output:}
\begin{verbatim}
a = 0; b = 3
f(x) = (x^2+x+1)/(x^2+1);  F(x) = integrate(f(x),x);  
show('f(x) = ',f(x), '  its antiderivative F(x) = ',F(x))

Int = integrate(f(x),x,a,b)
show('Integrate command gives = ', Int)
show('  also from Fundamental Theorem F(b)-F(a) = ',F(b)-F(a))
\end{verbatim}
\end{programcode}
The output is shown in Fig \ref{LogAntider2}
\begin{image}[15cm]
%\sidecaption[t]
\includegraphics[scale=.65]{/LogAntider2.png}
\caption{Antiderivatives of log}
\label{LogAntider2}       % Give a unique label
\end{image}


In some integrals like the one below, the need for the change of variable is not obvious, but simplifying the denominator helps.   SageMath can handle them just as easily.
$$ \int \frac{3x}{(x+2)^2}dx =  \int \frac{3x}{(x+2)^2}dx = [u=x+2 \Rightarrow \ du=dx] = $$
$$ \int \frac{3(u-2)}{u^2}du =  \int \frac{3}{u} - \frac{6}{u^2} =  3\ln(u)+ \frac{6}{u}= 3\ln(x+2)+ \frac{6}{x+2}$$
\begin{programcode}{SageMath Code and Output:}
\begin{verbatim}
a = -1; b = 4
f(x) = 3*x/(x+2)^2;  F(x) = integrate(f(x),x);  
show('f(x) = ',f(x), '  its antiderivative F(x) = ',F(x))

Int = integrate(f(x),x,a,b)
show('Integrate command gives = ', Int)
show('  also from Fundamental Theorem F(b)-F(a) = ',F(b)-F(a))
\end{verbatim}
\end{programcode}
The output is shown in Fig \ref{ChangeVar}
\begin{image}[15cm]
%\sidecaption[t]
\includegraphics[scale=.65]{/ChangeVar.png}
\caption{Change variable of integration}
\label{ChangeVar}       % Give a unique label
\end{image}


A classic integral with log substitution is handled just as well by SageMath.  
$$ \int \frac{1}{x\ln(x)}dx = \left[u=\ln(x) \ \Rightarrow \ du = \frac{1}{x}dx \right] = \int  \frac{1}{u}du = \ln|u| = \ln|\ln|x|| $$
$$ \int_e^{e^e} \frac{1}{x\ln(x)}dx = [u=\ln(x) \ \Rightarrow \ du = \frac{1}{x}dx] = \int_1^e  \frac{1}{u}du = \ln|u|_1^e = 1 $$
\begin{programcode}{SageMath Code and Output:}
\begin{verbatim}
a = e^1; b = e^e
f(x) = 1/(x*log(x));  F(x) = integrate(f(x),x);  
show('f(x) = ',f(x), '  its antiderivative F(x) = ',F(x))

Int = integrate(f(x),x,a,b)
show('Integrate command gives = ', Int)
show('  also from Fundamental Theorem F(b)-F(a) = ',F(b)-F(a))
\end{verbatim}
\end{programcode}
The output is shown in Fig \ref{LogLogInt}
\begin{image}[15cm]
%\sidecaption[t]
\includegraphics[scale=.65]{/LogLogInt.png}
\caption{Log-log integral}
\label{LogLogInt}       % Give a unique label
\end{image}


The integral below illustrates a simple change of variable as well as exponential integration with base-$2$ rather than base-$e$ that leads to extra $\ln(2)$ in the denominator.  $$\int 2^{x^2}xdx = [u=x^2 \ \ \Rightarrow \ \ du = 2xdx] = \int 2^u \frac{du}{2} = \frac{2^u}{2\ln(2)} = \frac{2^{x^2-1}}{\ln(2)}$$
\begin{programcode}{SageMath Code and Output:}
\begin{verbatim}
f(x) = 2^(x^2)*x;  F(x) = integrate(f(x),x)
show('f(x) = ',f(x),'  its antiderivative F(x) = ',F(x))
\end{verbatim}
\end{programcode}
The output is shown in Fig \ref{Int2x2}
\begin{image}[15cm]
%\sidecaption[t]
\includegraphics[scale=.65]{/Int2x2.png}
\caption{Integral of $2^{x^2}x$}
\label{Int2x2}       % Give a unique label
\end{image}






%%%%%%%%%%%%%%%%%%%%%%%%%%%%%%%%%%%%%%%%%%%%%%%%%%%%%%%%%%%%%%%%%%%%%%%%%%%
\section{Special Integration Techniques}

As we mentioned at the beginning of the chapter, SageMath just computes harder integrals requiring tedious analytical computations like a "black box".   Therefore, in the code below we don't repeat longer analytical solutions, just remark which type of analytical integration they refer to.   Sometimes SageMath (and other CAS systems) produces a form of the answer different from the textbook analytical solution, but equivalent after trigonometric substitutions and algebraic simplifications.  We begin with the SageMath cell below that computes several standard inverse trigonometric functions integrals with minimal substitutions. The third integral is usually written with inverse secant function, but the SageMath preferred $\arcsin(x)$.

\begin{programcode}{SageMath Code and Output:}
\begin{verbatim}
f(x) = 1/sqrt(36-4*x^2);  F(x) = integrate(f(x),x)
show('f(x) = ',f(x),'  antiderivative F(x) = ',F(x))
f(x) = 1/(9+4*x^2);  F(x) = integrate(f(x),x)
show('f(x) = ',f(x),'  antiderivative F(x) = ',F(x))
f(x) = 1/(x*sqrt(x^2-9));  F(x) = integrate(f(x),x)
show('f(x) = ',f(x),'  antiderivative F(x) = ',F(x))
f(x) = 1/(sqrt(e^(2*x)-9));  F(x) = integrate(f(x),x)
show('f(x) = ',f(x),'  antiderivative F(x) = ',F(x))
\end{verbatim}
\end{programcode}
The output is shown in Fig \ref{IntInverseTrig}
\begin{image}[15cm]
%\sidecaption[t]
\includegraphics[scale=.65]{/IntInverseTrig.png}
\caption{Inverse trig}
\label{IntInverseTrig}       % Give a unique label
\end{image}


With the analytical approach, the integrals below require a breakdown of the numerator into the term with $x$ that can be computed with a change of variable and a constant term, producing an inverse trigonometric function. SageMath gives standard answers.
\begin{programcode}{SageMath Code and Output:}
\begin{verbatim}
f(x) = (2*x+3)/sqrt(9-x^2);  F(x) = integrate(f(x),x)
show('f(x) = ',f(x),'  antiderivative F(x) = ',F(x))
f(x) = (2*x+3)/(x^2+9);  F(x) = integrate(f(x),x)
show('f(x) = ',f(x),'  antiderivative F(x) = ',F(x))
\end{verbatim}
\end{programcode}
The output is shown in Fig \ref{IntAsinAtan}
\begin{image}[15cm]
%\sidecaption[t]
\includegraphics[scale=.65]{/IntAsinAtan.png}
\caption{Inverse Trig}
\label{IntAsinAtan}       % Give a unique label
\end{image}

The integrals below require completion of squares and SageMath computes them easily.
\begin{programcode}{SageMath Code and Output:}
\begin{verbatim}
f(x) = 1/(2*x^2-12*x+26);  F(x) = integrate(f(x),x)
show('f(x) = ',f(x),'  antiderivative F(x) = ',F(x))
f(x) = 1/sqrt(16*x-4*x^2);  F(x) = integrate(f(x),x)
show('f(x) = ',f(x),'  antiderivative F(x) = ',F(x))
\end{verbatim}
\end{programcode}
The output is shown in Fig \ref{CompleteSquare}
\begin{image}[15cm]
%\sidecaption[t]
\includegraphics[scale=.65]{/CompleteSquare.png}
\caption{Complete Square}
\label{CompleteSquare}       % Give a unique label
\end{image}

In the next cell, SageMath performs several standard integrations by parts.
\begin{programcode}{SageMath Code and Output:}
\begin{verbatim}
f(x) = x*e^x;  F(x) = integrate(f(x),x)
show('f(x) = ',f(x),'  F(x) = ',F(x))
f(x) = x^2*e^x;  F(x) = integrate(f(x),x)
show('f(x) = ',f(x),'  F(x) = ',F(x))
f(x) = x^2*log(x);  F(x) = integrate(f(x),x)
show('f(x) = ',f(x),'  F(x) = ',F(x))
f(x) = x*cos(x);  F(x) = integrate(f(x),x)
show('f(x) = ',f(x),'  F(x) = ',F(x))
f(x) = x^2*cos(x);  F(x) = integrate(f(x),x)
show('f(x) = ',f(x),'  F(x) = ',F(x))
f(x) = arcsin(x);  F(x) = integrate(f(x),x)
show('f(x) = ',f(x),'  F(x) = ',F(x))
f(x) = e^x*cos(2*x);  F(x) = integrate(f(x),x)
show('f(x) = ',f(x),'  F(x) = ',F(x))
\end{verbatim}
\end{programcode}
The output is shown in Fig \ref{IntegrByParts}
\begin{image}[15cm]
%\sidecaption[t]
\includegraphics[scale=.65]{/IntegrByParts.png}
\caption{Integration by parts}
\label{IntegrByParts}       % Give a unique label
\end{image}


We illustrate several trigonometric integrals below that require substitutions of $u = \cos(x), \sin(x), \tan(x), \sec(x)$, etc... 
\begin{programcode}{SageMath Code and Output:}
\begin{verbatim}
f(x) = (cos(x))^3*(sin(x))^6;  F(x) = integrate(f(x),x)
show('f(x) = ',f(x),'  F(x) = ',F(x))
f(x) = (sin(x))^3/sqrt(cos(x));  F(x) = integrate(f(x),x)
show('f(x) = ',f(x),'  F(x) = ',F(x))
f(x) = (sin(x))^4;  F(x) = integrate(f(x),x)
show('f(x) = ',f(x),'  F(x) = ',F(x))
f(x) = (tan(x))^5/(sec(x))^(3/2);  F(x) = integrate(f(x),x)
show('f(x) = ',f(x),'  F(x) = ',F(x))
f(x) = (sec(x))^3;  F(x) = integrate(f(x),x)
show('f(x) = ',f(x),'  F(x) = ',F(x))
f(x) = sin(3*x)*cos(4*x);  F(x) = integrate(f(x),x)
show('f(x) = ',f(x),'  F(x) = ',F(x))
\end{verbatim}
\end{programcode}
The output is shown in Fig \ref{SinCosSubstitution}
\begin{image}[15cm]
%\sidecaption[t]
\includegraphics[scale=.65]{/SinCosSubstitution.png}
\caption{Sin/cos substitutions}
\label{SinCosSubstitution}       % Give a unique label
\end{image}



The integrals below require trigonometric substitutions $x=\arcsin(x), \arctan(x)$, etc...  These tend to be quite tedious analytically, but SageMath has no issues.
\begin{programcode}{SageMath Code and Output:}
\begin{verbatim}
f(x) = 1/(x^2*sqrt(16-x^2));  F(x) = integrate(f(x),x)
show('f(x) = ',f(x),'  F(x) = ',F(x))
f(x) = sqrt(4*x^2+9)/x^4;  F(x) = integrate(f(x),x)
show('f(x) = ',f(x),'  F(x) = ',F(x))
f(x) = 1/sqrt(4*x^2-9);  F(x) = integrate(f(x),x)
show('f(x) = ',f(x),'  F(x) = ',F(x))
f(x) = sqrt(x^2-9)/x;  F(x) = integrate(f(x),x)
show('f(x) = ',f(x),'  F(x) = ',F(x))
\end{verbatim}
\end{programcode}
The output is shown in Fig \ref{TrigSubstitution}
\begin{image}[15cm]
%\sidecaption[t]
\includegraphics[scale=.65]{/TrigSubstitution.png}
\caption{Inverse substitutions}
\label{TrigSubstitution}       % Give a unique label
\end{image}

It is particularly remarkable how well SageMath can deal with tedious partial fractions integrations shown below.   In the first integrant, the power of the numerator is higher than the denominator, so long division has to be done first, but \texttt{.partial\_fraction()} takes care of that too.

\begin{programcode}{SageMath Code and Output:}
\begin{verbatim}
f(x) = (2*x^4-3*x-1)/(2*x^3 - 5*x^2 + 3*x);  
F(x) = integrate(f(x),x)
show('f(x) = ',f(x),' = ', f(x).partial_fraction());
show('  F(x) = ',F(x))

f(x) = (4*x^2-2*x+1)/(x^3 - 4*x^2 + 4*x); 
F(x) = integrate(f(x),x)
show('f(x) = ',f(x),' = ', f(x).partial_fraction());
show('  F(x) = ',F(x))

f(x) = (3*x+1)/(x^4 + x^3 + 9*x^2 + 9*x);  
F(x) = integrate(f(x),x)
show('f(x) = ',f(x),' = ', f(x).partial_fraction());
show('  F(x) = ',F(x))

f(x) = (3*x^3+x^2-x-1)/(x^5 + 2*x^3 + x);  
F(x) = integrate(f(x),x)
show('f(x) = ',f(x),' = ', f(x).partial_fraction());
show('  F(x) = ',F(x))

f(x) = (1)/(x^4 + 2*x^2 + 1);  
F(x) = integrate(f(x),x)
show('f(x) = ',f(x),' = ', f(x).partial_fraction());
show('  F(x) = ',F(x))
\end{verbatim}
\end{programcode}
The output is shown in Fig \ref{PartialFractions}
\begin{image}[15cm]
%\sidecaption[t]
\includegraphics[scale=.65]{/PartialFractions.png}
\caption{Partial fractions}
\label{PartialFractions}       % Give a unique label
\end{image}



%%%%%%%%%%%%%%%%%%%%%%%%%%%%%%%%%%%%%%%%%%%%%%%%%%%%%%%%%%%%%%%%%%%%%%%%%%%
\section{Improper Integrals}
For the improper integrals, most of the time the computation consists of finding a regular integral on a finite interval and taking a limit (often requiring the application of L'Hopital's rule).
$$ \int_1^\infty \frac{1}{x^2}dx = lim_{b \rightarrow \infty} \int_1^b \frac{1}{x^2}dx = lim_{b \rightarrow\infty} \left( -\frac{1}{x} \right)_1^b =$$
$$ lim_{b \rightarrow \infty} \left( -\frac{1}{b} +1\right) =1$$
To illustrate the analytical approach, SageMath can be used to compute antiderivative first, substitute limits of integration, and then take the infinite limit, but direct integration will work as well.   The steps are quite algorithmic and are coded as a function below.  Note that we only perform the direct computation if the limit is finite; otherwise a rather verbose error is obtained.

\begin{programcode}{SageMath Code and Output:}
\begin{verbatim}
def improperintegral(x,f,a,b,c):
    F(x) = integrate(f(x),x); L = limit(F(b)-F(a),b=c)
    show('f(x) = ',f(x), ' F(x) = ', F(x))
    show(' after integration = ', F(b) - F(a))
    show(' limit at b = ', c, ' is = ', L)
    if ( L < infinity ): # only compute directly for finite L
        show('Direct integration = ', integrate(f(x),x,a,c))
    show(plot(f(x) ,(x,a,10*a), figsize = (4,3), gridlines=True))
    

f(x) = 1/x^2; a = 1; var('b'); c = infinity
improperintegral(x,f,a,b,c)
\end{verbatim}
\end{programcode}
The output is shown in Fig \ref{IntInfx2}
\begin{image}[15cm]
%\sidecaption[t]
\includegraphics[scale=.65]{/IntInfx2.png}
\caption{Improper integral of $\frac{1}{x^2}$}
\label{IntInfx2}       % Give a unique label
\end{image}



Analogously $$\int_1^\infty \frac{1}{\sqrt{x}}dx = lim_{b \rightarrow \infty} \int_1^b x^{-1/2}dx =$$
$$ lim_{b \rightarrow \infty} \left( \frac{x^{1/2}}{\frac{1}{2}} \right)_1^b = lim_{b \rightarrow \infty} \left( 2\sqrt{b}-2\right) =\infty$$
$$\int_1^\infty \frac{1}{x}dx = lim_{b \rightarrow \infty} \int_1^b  \frac{1}{x} dx =$$
$$ lim_{b \rightarrow \infty} \left( \ln(x) \right)_{1}^b = lim_{b \rightarrow \infty} (\ln(b)-ln(1))  = \infty$$

\begin{programcode}{SageMath Code and Output:}
\begin{verbatim}
f(x) = 1/sqrt(x); a = 1; var('b'); c = infinity
improperintegral(x,f,a,b,c)
f(x) = 1/x; a = 1; var('b'); c = infinity
improperintegral(x,f,a,b,c)
\end{verbatim}
\end{programcode}
The output is shown in Fig \ref{IntInfsqrt}
\begin{image}[15cm]
%\sidecaption[t]
\includegraphics[scale=.65]{/IntInfsqrt.png}
\caption{Improper integral of $\frac{1}{\sqrt{x}}$}
\label{IntInfsqrt}       % Give a unique label
\end{image}


The indefinite integrals due to the function going to infinity at one of the limits can be handled similarly
$$ \int_0^1 \frac{1}{x^2}dx = lim_{b \rightarrow 0} \int_b^1 \frac{1}{x^2}dx = lim_{b \rightarrow 0} \left( -\frac{1}{x} \right)_b^1 =$$
$$ lim_{b \rightarrow 0} \left(-1+\frac{1}{b}\right) =\infty$$
$$\int_0^1 \frac{1}{\sqrt{x}}dx = lim_{b \rightarrow 0} \int_b^1 x^{-1/2}dx = lim_{b \rightarrow 0} \left( \frac{x^{1/2}}{\frac{1}{2}} \right)_b^1 =$$
$$ lim_{b \rightarrow 0} \left( 2-2\sqrt{b}\right) =2$$

\begin{programcode}{SageMath Code and Output:}
\begin{verbatim}
reset(); var('b'); c = 0
f(x) = 1/x^2;  
F(x) = integrate(f(x),x); L = limit(F(b)-F(1),b=c)
show('f(x) = ',f(x), ' F(x) = ', F(x))
show(' after integration = ', F(1) - F(b))
show(' its limit at b = ', c, ' is = ', limit(F(1)-F(b),b=c))
p1 = plot(f(x),(x,0,1),ymax = 20,legend_label = '1/x^2')

f(x) = 1/sqrt(x);  
F(x) = integrate(f(x),x); L = limit(F(b)-F(1),b=c)
show('f(x) = ',f(x), ' F(x) = ', F(x))
show(' after integration = ', F(1) - F(b))
show(' its limit at b = ', c, ' is = ', limit(F(1)-F(b),b=c))
show(' Direct integration = ', integrate(f(x),x,0,1))
p2 = plot(f(x),(x,0,1),ymax = 20,legend_label = '1/sqrt(x)')

G = graphics_array([[p1, p2]])
G.show(figsize=[6, 3])
\end{verbatim}
\end{programcode}
The output is shown in Fig \ref{ImpoperX2sqrt}
\begin{image}[15cm]
%\sidecaption[t]
\includegraphics[scale=.65]{/ImpoperX2sqrt.png}
\caption{Improper integral of $\frac{1}{x^2}$ and $\frac{1}{\sqrt{x}} 0 to 1 limit$}
\label{ImpoperX2sqrt}       % Give a unique label
\end{image}


The SageMath cell below applies the limit twice to compute this doubly improper integral.
$$ \int_0^\infty \frac{1}{\sqrt{x}(x+1)}dx = lim_{a \rightarrow 0} lim_{b \rightarrow \infty}   \int_a^b \frac{1}{\sqrt{x}(x+1)} dx=$$
$$[u = \sqrt{x} \ \ \ du = \frac{1}{2\sqrt{x}}dx]=  lim_{a \rightarrow 0} lim_{b \rightarrow \infty}   \int_{\sqrt{a}}^{\sqrt{b}} \frac{1}{u^2+1} 2du =$$
$$ lim_{a \rightarrow 0} lim_{b \rightarrow \infty} \left( arctan(\sqrt{b}) -arctan(\sqrt{a}) \right) = \frac{\pi}{2}$$

\begin{programcode}{SageMath Code and Output:}
\begin{verbatim}
reset(); var('a,b')
f(x) = 1/(sqrt(x)*(x+1));  F(x) = integrate(f(x),x)
show('f(x) = ',f(x))
show(' after integration = ', F(b) - F(a))
show(' its double limit = ', \
     limit(limit(F(b) - F(a),a=0),b=infinity))
show('Direct integration = ', integrate(f(x),x,0,infinity))
plot(f(x),(x,0,10),ymax = 20,figsize = (4,3))
\end{verbatim}
\end{programcode}
The output is shown in Fig \ref{ImproperDouble}
\begin{image}[15cm]
%\sidecaption[t]
\includegraphics[scale=.65]{/ImproperDouble.png}
\caption{Double limits}
\label{ImproperDouble}       % Give a unique label
\end{image}


SageMath computes improper integrals requiring integration by parts just as easily.
\begin{programcode}{SageMath Code and Output:}
\begin{verbatim}
reset(); var('b'); c = infinity
f(x) = (x-4)*e^(-x);  F(x) = integrate(f(x),x)
show('f(x) = ',f(x), ' after integration = ', F(b) - F(1))
show(' its limit at b = ', c, ' is = ', limit(F(b)-F(1),b=c))
show('Direct integration = ', integrate(f(x),x,1,infinity))

reset(); var('b'); c = 0
f(x) = log(x);  F(x) = integrate(f(x),x)
show('f(x) = ',f(x), ' after integration = ', F(1) - F(b))
show(' its limit at b = ', c, ' is = ', limit(F(1)-F(b),b=c))
show('Direct integration = ', integrate(f(x),x,0,1))
\end{verbatim}
\end{programcode}
The output is shown in Fig \ref{ImproperbyParts}
\begin{image}[15cm]
%\sidecaption[t]
\includegraphics[scale=.65]{/ImproperbyParts.png}
\caption{Integration by parts, improper}
\label{ImproperbyParts}       % Give a unique label
\end{image}


Finally, SageMath can deal with the integral of $e^{-x^2}$ where antiderivative is a special function.
\begin{programcode}{SageMath Code and Output:}
\begin{verbatim}
reset(); var('b'); c = infinity
f(x) = e^(-x^2);  F(x) = integrate(f(x),x)
show('f(x) = ',f(x), ' after integration = ', F(b) - F(0))
show(' its limit at b = ', c, ' is = ', limit(F(b)-F(0),b=c))
show('Direct integration = ', integrate(f(x),x,0,infinity))
plot(f(x),(x,-5,5), figsize = (4,3))
\end{verbatim}
\end{programcode}
The output is shown in Fig \ref{ImproperEx2}
\begin{image}[15cm]
%\sidecaption[t]
\includegraphics[scale=.65]{/ImproperEx2.png}
\caption{Improper integral of $e^{-x^2}$}
\label{ImproperEx2}       % Give a unique label
\end{image}













%%%%%%%%%%%%%%%%%%%%%%%%%%%%%%%%%%%%%%%%%%%%%%%%%%%%%%%%%%%%%%%%%%%%%%%%%%%
\section{Applications of Integration}

\subsection{Area Between Curves} %%%%%%%%%%%%%%%%%%%%%%%%%%%%%%%%%%%%%%%%%%%%%%%%

Let's begin the discussion of the applications of integration by finding the areas between the two curves.  In the example below, we have to first plot the curves and find the intersection points, and then, in the next SageMath cell, we compute the actual integral equal to the area between the functions. 


\begin{programcode}{SageMath Code and Output:}
\begin{verbatim}
reset();   f(x) = 6-x^2;  g(x) = -x;
sol = solve(f(x)==g(x),x); 
show('Intersection points are: ',sol)
a = -5; b = 5; c = -10; d = 10; # plot bounds
plot1 = plot(f(x), (x,a,b), ymin=c, ymax=d, \
 legend_label= "f(x)",color='red', figsize = (4,3), gridlines=True)
plot2 = plot(g(x), (x,a,b), ymin=c, ymax=d, \
 legend_label= "g(x)",color='blue',figsize = (4,3), gridlines=True)
plot_all = plot1 + plot2; plot_all.show()
\end{verbatim}
\end{programcode}
The output is shown in Fig \ref{AreaBetwgraph1}
\begin{image}[15cm]
%\sidecaption[t]
\includegraphics[scale=.65]{/AreaBetwgraph1.png}
\caption{Plot curves}
\label{AreaBetwgraph1}       % Give a unique label
\end{image}


\begin{programcode}{SageMath Code and Output:}
\begin{verbatim}
differ(x) = f(x) - g(x);  
F(x) = integrate(differ(x),x);  L = -2;  R = 3
show('differ(x) = ',differ(x))
show('  its antiderivative F(x) = ',F(x))
Int = integrate(differ(x),x,L,R)
show('Integrate command gives = ', Int)
show('  also F(R)-F(L) = ',F(R)-F(L))
\end{verbatim}
\end{programcode}
The output is shown in Fig \ref{AreaBetw1}
\begin{image}[15cm]
%\sidecaption[t]
\includegraphics[scale=.65]{/AreaBetw1.png}
\caption{Area between curves}
\label{AreaBetw1}       % Give a unique label
\end{image}

In the simpler example above, we can illustrate the analytical solution of the problem as well:
$$ f(x) = 6-x^2 \ \ \ \ \ g(x) = -x \ \ \ \ \text{intersections: } \ \ \ \  6-x^2 = -x \Rightarrow $$
$$ x^2-x-6=(x-3)(x+2)=0 \Rightarrow \ \ \ \ x = -2,3 $$
$$ A = \int_a^b f(x) - g(x) dx = \int_{-2}^3 ((6-x^2)-(-x))dx =$$
$$  \int_{-2}^3 (-x^2+x+6)dx =\left( -\frac{x^3}{3} +\frac{x^2}{2} +6x \right)_{-2}^3 = \frac{125}{6}$$



The SageMath cell below illustrates the case when there are more than two curve intersections, so the area is the sum of areas in between each intersection, and which curve is above or below has to be specified individually on each interval.

\begin{programcode}{SageMath Code and Output:}
\begin{verbatim}
reset();   f(x) = 2*x^3+x^2-2*x;  g(x) = x^3+2*x^2;
sol = solve(f(x)==g(x),x); show(sol)
a = -2; b = 3; c = -10; d = 20; # plot bounds
plot1 = plot(f(x), (x,a,b), ymin=c, ymax=d,  \
    legend_label= "f(x)",color='red', figsize = (4,3), gridlines=True)
plot2 = plot(g(x), (x,a,b), ymin=c, ymax=d, \
    legend_label= "g(x)",color='blue', figsize = (4,3), gridlines=True)
plot_all = plot1 + plot2; plot_all.show()
\end{verbatim}
\end{programcode}
The output is shown in Fig \ref{AreaBetwTwograph}
\begin{image}[15cm]
%\sidecaption[t]
\includegraphics[scale=.65]{/AreaBetwTwograph.png}
\caption{Two intersections}
\label{AreaBetwTwograph}       % Give a unique label
\end{image}

Based on the above output. we can compute:

\begin{programcode}{SageMath Code and Output:}
\begin{verbatim}
differ1(x) = f(x) - g(x);  
F1(x) = integrate(differ1(x),x);  L = -1;  R = 0
show('differ1(x) = ',differ1(x))
show('  its antiderivative F1(x) = ',F1(x))
Int1 = integrate(differ1(x),x,L,R)
show('Integrate command gives = ', Int1)
show('  also F1(R)-F1(L) = ',F1(R)-F1(L))

differ2(x) = g(x) - f(x);  
F2(x) = integrate(differ2(x),x);  L = 0;  R = 2
show('differ2(x) = ',differ2(x))
show('  its antiderivative F2(x) = ',F2(x))
Int2 = integrate(differ2(x),x,L,R)
show('Integrate command gives = ', Int2)
show('  also F2(R)-F2(L) = ',F2(R)-F2(L))

show('\nTotal area between curves = Int1 + Int2 = ', Int1 + Int2)
\end{verbatim}
\end{programcode}
The output is shown in Fig \ref{AreaBetwTwoCompute}
\begin{image}[15cm]
%\sidecaption[t]
\includegraphics[scale=.65]{/AreaBetwTwoCompute.png}
\caption{Area between two intersections}
\label{AreaBetwTwoCompute}       % Give a unique label
\end{image}


If the functions of $y$ are given, then the integration must be done in $y$ as the next two SageMath cells show.  Note that in this case \texttt{implicit\_plot()} is necessary to plot implicit equations.

\begin{programcode}{SageMath Code and Output:}
\begin{verbatim}
reset();  var('y');   f(y) = 2-y^2;  g(y) = y;
sol = solve(f(y)==g(y),y); show(sol)
a = -5; b = 3; c = -3; d = 3; # plot bounds
plot1 = implicit_plot(x==f(y), (x,a,b), (y,c,d), \
        legend_label= "f",color='red', figsize = (4,3), gridlines=True)
plot2 = implicit_plot(x==g(y), (x,a,b), (y,c,d), \
        legend_label= "g",color='blue', figsize = (4,3), gridlines=True)
plot_all = plot1 + plot2; plot_all.show()
\end{verbatim}
\end{programcode}
The output is shown in Fig \ref{AreaBetwYgraph}
\begin{image}[15cm]
%\sidecaption[t]
\includegraphics[scale=.65]{/AreaBetwYgraph.png}
\caption{Plot functions of y}
\label{AreaBetwYgraph}       % Give a unique label
\end{image}

Based on the above output. we can see the bounds in $y$ and can compute as follows:

\begin{programcode}{SageMath Code and Output:}
\begin{verbatim}
differ(y) = f(y) - g(y);  
F(y) = integrate(differ(y),y);  L = -2;  R = 1
show('differ(y) = ',differ(y))
show('  its antiderivative F(y) = ',F(y))
Int = integrate(differ(y),y,L,R)
show('Integrate command gives = ', Int)
show('  also F(R)-F(L) = ',F(R)-F(L))
\end{verbatim}
\end{programcode}
The output is shown in Fig \ref{AreaBetwYcompute}
\begin{image}[15cm]
%\sidecaption[t]
\includegraphics[scale=.65]{/AreaBetwYcompute.png}
\caption{Area between function of y}
\label{AreaBetwYcompute}       % Give a unique label
\end{image}







\subsection{Volume: Disk and Shell Methods} %%%%%%%%%%%%%%%%%%%%%%%%%%%%%%%%%%%%%%%%%%%%%%%%

The next SageMath cell illustrates a classic problem of finding volume between two curves rotating around the $x$-axis via the disk or washer method. First, we find intersection points and plot a 2-D region that is rotated. Then the disk/washer circular cross sections are used to find the volume.
\begin{programcode}{SageMath Code and Output:}
\begin{verbatim}
reset();   f(x) = x;  g(x) = x^2;
sol = solve(f(x)==g(x),x); show(sol)
a = -1; b = 3; c = -1; d = 2; # plot bounds
plot1 = plot(f(x), (x,a,b), ymin=c, ymax=d,  
                legend_label= "f(x)", color='red',  
                figsize = (4,3), gridlines=True)
plot2 = plot(g(x), (x,a,b), ymin=c, ymax=d,
                legend_label= "g(x)", color='blue', 
                figsize = (4,3), gridlines=True)
plot_all = plot1 + plot2; plot_all.show()
\end{verbatim}
\end{programcode}
The output is shown in Fig \ref{VolRotatXgraph}
\begin{image}[15cm]
%\sidecaption[t]
\includegraphics[scale=.65]{/VolRotatXgraph.png}
\caption{Intersection of curves}
\label{VolRotatXgraph}       % Give a unique label
\end{image}
Based on the above output. we can see the bounds in $x$ and can compute volume as follows:
\begin{programcode}{SageMath Code and Output:}
\begin{verbatim}
Vp(x) = pi*(x^2-(x^2)^2);   V(x) = integrate(Vp(x),x)
show('Vp = ', Vp(x), ' its antiderivative = ', V(x))
show('Volume = ',V(1)-V(0)) 
\end{verbatim}
\end{programcode}
The output is shown in Fig \ref{VolRotatX}
\begin{image}[15cm]
%\sidecaption[t]
\includegraphics[scale=.65]{/VolRotatX.png}
\caption{Volume between curves}
\label{VolRotatX}       % Give a unique label
\end{image}
Once again, in the simpler example above, we can illustrate the analytical solution to the problem as well:
$$ R(x) =x \ \ \ \ \ r(x) = x^2 \ \ \ \ \text{intersections: } \ \ \ \  x = x^2 \Rightarrow $$
$$ x(1-x)=0 \Rightarrow \ \ \ \ x = 0,1 $$
$$ V = \pi \int_a^b \pi(R(x)^2 - r(x)^2) dx =\pi \int_0^1 (x^2-x^4)dx = \pi
 \left( \frac{x^3}{3} - \frac{x^5}{5} \right)_0^1 = \frac{2\pi}{15}$$

Next, the same region above is rotated around the horizontal line $y=3$.   The outside radius becomes $3-x^2$, while the inner radius is now $3-x$.   Therefore
$$ V = \pi \int_a^b \pi(R(x)^2 - r(x)^2) dx =\pi \int_0^1 ((3-x^2)^2-(3-x)^2)dx $$

\begin{programcode}{SageMath Code and Output:}
\begin{verbatim}
Vp(x) = pi*((3-x^2)^2-(3-x)^2);   V(x) = integrate(Vp(x),x)
show('Vp = ', Vp(x), ' or expanded = ', expand(Vp(x)))
show('  its antiderivative = ', V(x))
show('Volume = ',V(1)-V(0)) 
\end{verbatim}
\end{programcode}
The output is shown in Fig \ref{VolRotatY3}
\begin{image}[15cm]
%\sidecaption[t]
\includegraphics[scale=.65]{/VolRotatY3.png}
\caption{Rotating about $y=3$}
\label{VolRotatY3}       % Give a unique label
\end{image}

The same region is now rotated about vertical line $x = -2$.  In this case, the variable of integration is $y$ and fuctions have to be rewritten as: 
$$ y=f(x)=x \ \ \ \text{therefore} \ \ \ x = y \ \ \ \ \text{and}  \ \ \ y=f(x)=x^2 \ \ \ \text{therefore} \ \ \ x = \sqrt{y}$$
The outer radius is $\sqrt{y}+2$ and the inner one is $y+2$.


\begin{programcode}{SageMath Code and Output:}
\begin{verbatim}
Vp(y) = pi*((2+sqrt(y))^2-(2+y)^2);   
V(y) = integrate(Vp(y),y)
show('Vp = ', Vp(y), ' or expanded = ', expand(Vp(y)))
show('  its antiderivative = ', V(y))
show('Volume = ',V(1)-V(0)) 
\end{verbatim}
\end{programcode}
The output is shown in Fig \ref{VolRotatX2}
\begin{image}[15cm]
%\sidecaption[t]
\includegraphics[scale=.65]{/VolRotatX2.png}
\caption{Rotating about $x=2$}
\label{VolRotatX2}       % Give a unique label
\end{image}


As a practical application of the disk/washer method consider a cylindrical hole of radius 6 inches drilled in a sphere of radius 10 inches.   It can be thought of as a rotation around $x$-axis of a part of the upper circle $y= \sqrt{100-x^2}$ above $y=6$, To find the $x$-limits, let's find the intersections of the circle of radius 10 and the line $y=6$:  $\ \ \ \ x^2 + 6^2 = 10^2 \ \ \Rightarrow \ \ x^2 = 64 \ \ \Rightarrow \ \ x=\pm 8$.
\begin{programcode}{SageMath Code and Output:}
\begin{verbatim}
reset();   f(x) = sqrt(100-x^2);  g(x) = 6;
sol = solve(f(x)==g(x),x); show(sol)
a = -10; b = 10; c = -3; d = 10; # plot bounds
plot1 = plot(f(x), (x,a,b), ymin=c, ymax=d, figsize = (4,3), \
   legend_label= "f(x)",color='red', gridlines=True)
plot2 = plot(g(x), (x,a,b), ymin=c, ymax=d, \
   legend_label= "g(x)",color='blue', gridlines=True)
plot_all = plot1 + plot2; plot_all.show()
\end{verbatim}
\end{programcode}
The output is shown in Fig \ref{VolDrillSphere}
\begin{image}[15cm]
%\sidecaption[t]
\includegraphics[scale=.65]{/VolDrillSphere.png}
\caption{Circle and a line }
\label{VolDrillSphere}       % Give a unique label
\end{image}
Based on the above output. we can see the bounds in $x$ and can compute volume as follows:
\begin{programcode}{SageMath Code and Output:}
\begin{verbatim}
Vp(x) = pi*((sqrt(100-x^2))^2-6^2);   V(x) = integrate(Vp(x),x)
show('Vp = ', Vp(x), ' its antiderivative = ', V(x))
show('Volume = ',V(8)-V(-8)) 
\end{verbatim}
\end{programcode}
The output is shown in Fig \ref{VolSphereComput}
\begin{image}[15cm]
%\sidecaption[t]
\includegraphics[scale=.65]{/VolSphereComput.png}
\caption{Volume of drilled sphere}
\label{VolSphereComput}       % Give a unique label
\end{image}


Let's also consider a more algebraically tedious example rotating intersecting parabolas about the x-axis between $x=0$ and $x=3$.  Note that they intersect at $x=2$, so we need to take two separate volume integrals depending on which curves provides outer and inner radii.
\begin{programcode}{SageMath Code and Output:}
\begin{verbatim}
reset();   f(x) = x^2+1;  g(x) = -x^2+2*x+5;
sol = solve(f(x)==g(x),x); show(sol)
a = -2; b = 3; c = -1; d = 10; # plot bounds
plot1 = plot(f(x), (x,a,b), ymin=c, ymax=d,  
                legend_label= "f(x)", color='red',  
                figsize = (4,3), gridlines=True)
plot2 = plot(g(x), (x,a,b), ymin=c, ymax=d,
                legend_label= "g(x)", color='blue', 
                figsize = (4,3), gridlines=True)
plot_all = plot1 + plot2; plot_all.show()
\end{verbatim}
\end{programcode}
The output is shown in Fig \ref{MiddleIntersect}
\begin{image}[15cm]
%\sidecaption[t]
\includegraphics[scale=.65]{/MiddleIntersect.png}
\caption{Intersection of curves}
\label{MiddleIntersect}       % Give a unique label
\end{image}

Based on the above output. compute separate volumes and add them to get the final answer:
\begin{programcode}{SageMath Code and Output:}
\begin{verbatim}
Vp1(x) = pi*(g(x)^2-f(x)^2);   V1(x) = integrate(Vp1(x),x)
show('Vp1 = ', Vp1(x), ' its antiderivative = ', V1(x))
Vol1 = V1(2)-V1(0)
show('Vol1 = ',Vol1) 

Vp2(x) = pi*(f(x)^2-g(x)^2);   V2(x) = integrate(Vp2(x),x)
show('Vp2 = ', Vp2(x), ' its antiderivative = ', V2(x))
Vol2 = V2(3)-V2(2)
show('Vol2 = ',Vol2) 

show('Total volume = ', Vol1+Vol2)
\end{verbatim}
\end{programcode}
The output is shown in Fig \ref{VolRotatIntersect}
\begin{image}[15cm]
%\sidecaption[t]
\includegraphics[scale=.65]{/VolRotatIntersect.png}
\caption{Rotation of intersecting curves}
\label{VolRotatIntersect}       % Give a unique label
\end{image}



As a first illustration of the shell method, the SageMath cell below finds the volume obtained by rotating about the $y$-axis the region bounded by $y = 3x^2-x^3$ and $y=0$.   It cannot be solved with the disk/washer method because it would require solving a cubic equation for the radius bounds.

\begin{programcode}{SageMath Code and Output:}
\begin{verbatim}
reset();   f(x) = 3*x^2-x^3;  g(x) = 0;
sol = solve(f(x)==g(x),x); show(sol)
a = -0.5; b = 4; c = -1; d = 4; # plot bounds
plot1 = plot(f(x), (x,a,b), ymin=c, ymax=d, figsize = (4,3), \
  legend_label= "f(x)",color='red', gridlines=True)
plot2 = plot(g(x), (x,a,b), ymin=c, ymax=d, \
  legend_label= "g(x)",color='blue', gridlines=True)
plot_all = plot1 + plot2; plot_all.show()
\end{verbatim}
\end{programcode}
The output is shown in Fig \ref{VolShell1graph}
\begin{image}[15cm]
%\sidecaption[t]
\includegraphics[scale=.65]{/VolShell1graph.png}
\caption{Function shells }
\label{VolShell1graph}       % Give a unique label
\end{image}
Based on the above output. we can see the bounds in $x$ and can compute volume as follows:
\begin{programcode}{SageMath Code and Output:}
\begin{verbatim}
r = x; # radius of the shell
h = f(x); # height of the shell
Vp(x) = 2*pi*r*h; # element of the shell volume
V(x) = integrate(Vp(x),x)
show('Vp = ', Vp(x))
show('  its antiderivative = ', V(x))
show('Volume = ',V(3)-V(0)) 
\end{verbatim}
\end{programcode}
The output is shown in Fig \ref{VolShell1}
\begin{image}[15cm]
%\sidecaption[t]
\includegraphics[scale=.65]{/VolShell1.png}
\caption{Shell method}
\label{VolShell1}       % Give a unique label
\end{image}

If the region above is rotated around the vertical $x=4$ axis instead of the vertical $y$-axis, we only have to change the radius of the shell to$r = 4-x$, the heights are still function values $h=f(x)$.
\begin{programcode}{SageMath Code and Output:}
\begin{verbatim}
r = 4-x; h = f(x);  
Vp(x) = 2*pi*r*h;   V(x) = integrate(Vp(x),x)
show('Vp = ', Vp(x))
show('  its antiderivative = ', V(x))
show('Volume = ',V(3)-V(0)) 
\end{verbatim}
\end{programcode}
The output is shown in Fig \ref{VolShell2}
\begin{image}[15cm]
%\sidecaption[t]
\includegraphics[scale=.65]{/VolShell2.png}
\caption{Volume about  $x = 4$ axis}
\label{VolShell2}       % Give a unique label
\end{image}


The example below illustrates the shell method for rotating the region between $y=x^2$ and $y=x^3$ around $y$-axis.
\begin{programcode}{SageMath Code and Output:}
\begin{verbatim}
reset();   f(x) = x^2;  g(x) = x^3;
sol = solve(f(x)==g(x),x); show(sol)
a = -0.5; b = 2; c = -1; d = 2; # plot bounds
plot1 = plot(f(x), (x,a,b), ymin=c, ymax=d, figsize = (4,3), \
   legend_label= "f(x)",color='red', gridlines=True)
plot2 = plot(g(x), (x,a,b), ymin=c, ymax=d, \
   legend_label= "g(x)",color='blue', gridlines=True)
plot_all = plot1 + plot2; plot_all.show()
\end{verbatim}
\end{programcode}
The output is shown in Fig \ref{VolShell3graph}
\begin{image}[15cm]
%\sidecaption[t]
\includegraphics[scale=.65]{/VolShell3graph.png}
\caption{Functions $y=x^2$ and $y=x^3$}
\label{VolShell3graph}       % Give a unique label
\end{image}
Based on the above output. we can see the bounds in $x$ and can compute volume as follows:
\begin{programcode}{SageMath Code and Output:}
\begin{verbatim}
r = x; # radius of the shell
h = f(x)-g(x); # height of the shell
Vp(x) = 2*pi*r*h; # element of the shell volume
V(x) = integrate(Vp(x),x)
show('Vp = ', Vp(x),' expand(Vp) = ',expand(Vp(x)))
show(' its antiderivative = ', V(x))
show('Volume = ',V(1)-V(0)) 
\end{verbatim}
\end{programcode}
The output is shown in Fig \ref{VolShell3}
\begin{image}[15cm]
%\sidecaption[t]
\includegraphics[scale=.65]{/VolShell3.png}
\caption{Shell between $y=x^2$ and $y=x^3$}
\label{VolShell3}       % Give a unique label
\end{image}













\subsection{Arc Length and Surfaces of Revolution.} %%%%%%%%%%%%%%%%%%%%%%%%%%%%%%%%%%%%%%%%%%%%%%

The formula for arc length $s = \int_a^b \sqrt{1+\left( \frac{dy}{dx} \right)^2}dx$ often leads to integrals that are hard or impossible to compute analytically.   We first illustrate a few special cases where such calculation is possible.
For example:
$$ f(x) = \frac{x^3}{6}+\frac{1}{2x}  \Rightarrow \ \ \ \ \  f'(x) =  \frac{3x^2}{6}-\frac{1}{2x^2} =  \frac{1}{2} \left( x^2 - \frac{1}{x^2} \right)$$
$$ 1+ (f'(x))^2 = 1 +  \frac{1}{4} \left( x^2 - \frac{1}{x^2} \right)^2 =   \frac{1}{4} \left(4+x^4-2x^2 \frac{1}{x^2}+ \frac{1}{x^4} \right)=$$
$$\frac{1}{4} \left(x^4+2+ \frac{1}{x^4} \right) =\frac{1}{4} \left( x^2 + \frac{1}{x^2} \right)^2  $$
$$s = \int_a^b \sqrt{1+\left( \frac{dy}{dx} \right)^2}dx =  \int_1^3  \sqrt{\frac{1}{4} \left( x^2 + \frac{1}{x^2} \right)^2  }dx =  \int_1^3 \frac{1}{2} \left( x^2 + \frac{1}{x^2} \right)dx = $$
$$  \int_1^3 \frac{1}{2} \left( x^2 + x^{-2} \right)dx =\frac{1}{2} \left(  \frac{x^3}{3} - \frac{1}{x} \right)_1^3 =\frac{14}{3} $$
The calculation above lends itself well to an algorithm, so I coded it as a function below.

\begin{programcode}{SageMath Code and Output:}
\begin{verbatim}
def arc_length(x,f,a,b):
    assume(x>0);
    fp(x) = derivative(f(x),x)
    show('f(x) = ', f(x))
    show(' fp(x) = ', fp(x))
    under_root = 1 + (fp(x))^2;  
    under_factor = factor(1 + (fp(x))^2)
    show('under_root = ', under_root)
    show('  under_factor = ', under_factor)
    integrand = sqrt(under_factor);  
    F(x) = integrate(integrand,x)
    show('integrand = ', integrand.full_simplify())
    show(' antiderivative = ',F(x))
    Int = integrate(sqrt(1 + (fp(x))^2),x,a,b)
    show('Integrate command gives = ', Int)
    show('  also F(b)-F(a) = ',F(b)-F(a))
    show(plot(f(x),(x,a,b), figsize = (4,3), gridlines=True)) 

f(x) = x^3/6+1/(2*x);   a = 1; b = 3;
arc_length(x,f,a,b)
\end{verbatim}
\end{programcode}
The output is shown in Fig \ref{ArcLength1}
\begin{image}[15cm]
%\sidecaption[t]
\includegraphics[scale=.65]{/ArcLength1.png}
\caption{Arc length}
\label{ArcLength1}       % Give a unique label
\end{image}



Another such example is provided by:
$$ f(x) = 6x^{3/2}  \Rightarrow \ \ \ \ \  f'(x) = 6 \frac{3}{2}x^{1/2}= 9x^{1/2}$$
$$ 1+ (f'(x))^2 = 1 + 81x   $$
$$s =\int_a^b \sqrt{1+\left( \frac{dy}{dx} \right)^2}dx =  \int_0^1 \sqrt{1 + 81x}dx =  [u= 1+81x  \ \ \ du=81dx] =  $$
$$  \int_1^{82} \frac{1}{81} \sqrt{u}du =\frac{1}{81} \left(  \frac{u^{3/2}}{{3/2}}\right)_1^{82} =\frac{2}{243} \left(  82^{3/2} -1\right)$$


\begin{programcode}{SageMath Code and Output:}
\begin{verbatim}
f(x) = 6*x^(3/2);   a = 0; b = 1;
arc_length(x,f,a,b)
\end{verbatim}
\end{programcode}
The output is shown in Fig \ref{ArcLength2}
\begin{image}[15cm]
%\sidecaption[t]
\includegraphics[scale=.65]{/ArcLength2.png}
\caption{Arc length of 3/2 curve}
\label{ArcLength2}       % Give a unique label
\end{image}




As a practical example of arc length let's compute the length of a wire hanging between the polls.  It is called a catenary.
$$ y=f(x)=50(e^{x/100}+e^{-x/100})  \Rightarrow \ \ \ \ y' = \frac{1}{2}(e^{x/100}-e^{-x/100})  $$
$$ 1+ (y'))^2 = 1 +  \frac{1}{4} \left( e^{x/100}-e^{-x/100}\right)^2 = $$
$$\frac{1}{4} \left(4+e^{x/50}- 2e^{x/100}e^{-x/100}+ e^{-x/50} \right)= \frac{1}{4}(e^{x/100}+e^{-x/100})^2 $$
$$s = \int_a^b \sqrt{1+\left( \frac{dy}{dx} \right)^2}dx =  \int_{-100}^{100}  \sqrt{ \frac{1}{4}(e^{x/100}+e^{-x/100})^2 }dx $$
$$  \int_{-100}^{100} \frac{1}{2}(e^{x/100}+e^{-x/100})dx= \frac{1}{2}100 \left( e^{x/100}-e^{-x/100} \right) _{-100}^{100}=100(e-1/e) $$
The simplifications of SageMath below are a bit different, but the answer is the same.

\begin{programcode}{SageMath Code and Output:}
\begin{verbatim}
f(x) = 50*(e^(x/100)+e^(-x/100));  a = -100; b = 100;
arc_length(x,f,a,b)   
\end{verbatim}
\end{programcode}
The output is shown in Fig \ref{ArcLength3}
\begin{image}[15cm]
%\sidecaption[t]
\includegraphics[scale=.65]{/ArcLength3.png}
\caption{Length of wire}
\label{ArcLength3}       % Give a unique label
\end{image}



Note that the choice of 50 in front of the function is crucial to obtain a perfect square above to be able to integrate analytically.   The SageMath cell below computes a generic case numerically via the $numerical\_integral()$ function.   The imaginary part shown with `I` is negligible `e-12`.

\begin{programcode}{SageMath Code and Output:}
\begin{verbatim}
reset(); a = -100; b = 100;
f(x) = 60*(e^(x/100)+e^(-x/100));   
fp(x) = derivative(f(x),x)
show('f(x) = ', f(x))
show(' fp(x) = ', fp(x))
under_root = 1 + (fp(x))^2;  
under_factor = factor(1 + (fp(x))^2)
show('under_root = ', under_root)
show('  under_factor = ', under_factor)
Int = numerical_integral(sqrt(1 + (fp(x))^2),a,b)
show('Integrate command gives = ', N(Int))
plot(f(x),(x,a,b), ymin = 0, ymax = 200,
     figsize = (4,3), gridlines=True)
\end{verbatim}
\end{programcode}
The output is shown in Fig \ref{ArcLength4}
\begin{image}[15cm]
%\sidecaption[t]
\includegraphics[scale=.65]{/ArcLength4.png}
\caption{Length of wire 60}
\label{ArcLength4}       % Give a unique label
\end{image}


Somewhat similar integrals need to be dealt with in computations of areas of revolutions.  For example consider the area of the surface formed by revolving the graph of $y=f(x) = x^3$ on the interval $[0, 1]$ about the $x$-axis. 
$$ y' = 3x^2 \Rightarrow  \ \ \ \ \ \ S = 2 \pi \int_a^b r \sqrt{1+(f')^2} dx = 2 \pi \int_0^1 x^3 \sqrt{1+(3x^2)^2} dx  = $$
$$ [u = 1+9x^4 \ \ \  du= 36x^3dx] =  2 \pi \int_1^{10} \sqrt{u} \frac{du}{36}  = \left(\frac{\pi}{18} \frac{u^{\frac{3}{2}}}{\frac{3}{2}}\right)_1^{10}\approx 3.56 $$ 
In the SageMath function below we code the computation.  Note that depending on whether the rotation is around $x$ or $y$ axis, the radius of rotation changes accordingly.   We also use \texttt{revolution\_plot3d()} function to show 3-d surface of revolution, which can be rotated by a mouse to get a better view.

\begin{programcode}{SageMath Code and Output:}
\begin{verbatim}
def surface_area(x,f,a,b,around):
    assume(x>0);  
    if around =='x':
        r = f(x); # radius of rotation
    elif around =='y':
        r = x
    else:
        print('wrong axis of rotation')
        return
    
    fp(x) = derivative(f(x),x)
    show('f(x) = ', f(x))
    show(' fp(x) = ', fp(x))
    under_root = 1 + (fp(x))^2;  
    under_factor = factor(1 + (fp(x))^2)
    show('under_root = ', under_root)
    show('  under_factor = ', under_factor)
    integrand = 2*pi*r*sqrt(under_factor);  
    F(x) = integrate(integrand,x)
    show('integrand = ', integrand.full_simplify())
    show(' antiderivative = ',F(x))
    Int = integrate(2*pi*r*sqrt(1 + (fp(x))^2),x,a,b)
    show('Integrate command gives = ', N(Int))
    show('  also F(b)-F(a) = ',N(F(b)-F(a)))
    show(plot(f(x),(x,a,b), figsize = (4,3), gridlines=True)) 

f(x) = x^3; a = 0; b = 1;
show(revolution_plot3d((x,f(x),0), (x, a, b), (0, 2*pi), parallel_axis='x'))
surface_area(x,f,a,b,'x')
\end{verbatim}
\end{programcode}
The output is shown in Fig \ref{AreaRevolut1}
\begin{image}[15cm]
%\sidecaption[t]
\includegraphics[scale=.65]{/AreaRevolut1.png}
\caption{Revolve $x^3$ around $x$}
\label{AreaRevolut1}       % Give a unique label
\end{image}

%The rotation of $y=f(x)=\sqrt{4-x^2}$ from $-1$ to $1$ around $x$-axis results in cancellation of x-dependence in the integral:
%\begin{programcode}{SageMath Code and Output:}
%\begin{verbatim}
%f(x) = sqrt(4-x^2); a = -1; b = 1;
%surface_area(x,f,a,b,'x')
%\end{verbatim}
%\end{programcode}
%The output is shown in Fig \ref{AreaRevolut2}
%\begin{image}[15cm]
%%\sidecaption[t]
%\includegraphics[scale=.65]{/AreaRevolut2.png}
%\caption{Area of revolution of $y=f(x)=\sqrt{4-x^2}$}
%\label{AreaRevolut2}       % Give a unique label
%\end{image}

Finally, in the SageMath cell below, we compute the area of the figure obtained by rotation of $y=f(x)=x^2$ around $y$-axis.  Note how the radius of rotation changes to $r=x$ and the rotation axis must be adjusted.
\begin{programcode}{SageMath Code and Output:}
\begin{verbatim}
f(x) = x^2; a = 0; b = 1;
var('y')
show(revolution_plot3d((sqrt(y),y,0), (y, a, b), (0, 2*pi), parallel_axis='y'))
surface_area(x,f,a,b,'y')

\end{verbatim}
\end{programcode}
The output is shown in Fig \ref{AreaRevolut3}
\begin{image}[15cm]
%\sidecaption[t]
\includegraphics[scale=.65]{/AreaRevolut3.png}
\caption{Revolve  $x^2$ around $y$}
\label{AreaRevolut3}       % Give a unique label
\end{image}


\end{document}